\subsection{Changing the Basis: New Coordinates for the Same Vector}

Up to this point every vector has been described with respect to the horizontal
and vertical coordinate directions. That choice was natural because the
Cartesian axes had already been constructed, but it was still only a choice.
The geometric vector itself does not contain a preferred pair of coordinate
axes.

We now ask a new question:

\begin{center}
\emph{What numbers describe the same vector when we choose two different
coordinate directions?}
\end{center}

This is the problem of changing basis. The vector will remain fixed. Only the
language used to describe it will change.

\subsubsection*{A basis is a pair of independent directions}

\begin{definitionbox}
Two nonparallel vectors $\mathbf b_1$ and $\mathbf b_2$ form a basis of the
plane. This means that every vector $\mathbf u$ can be written in one and only
one way as
\[
\boxed{\mathbf u=c_1\mathbf b_1+c_2\mathbf b_2}.
\]
The numbers $c_1$ and $c_2$ are called the coordinates of $\mathbf u$ in the
basis
\[
\mathcal B=(\mathbf b_1,\mathbf b_2).
\]
We write
\[
[\mathbf u]_{\mathcal B}=(c_1,c_2).
\]
\end{definitionbox}

The basis vectors need not be perpendicular and need not have length one. The
only forbidden situation is that they are parallel, because then both vectors
point along the same line and cannot describe arbitrary displacements in the
plane.

Geometrically, the equation
\[
\mathbf u=c_1\mathbf b_1+c_2\mathbf b_2
\]
means the following. Starting at the origin, move first by
$c_1\mathbf b_1$, then by $c_2\mathbf b_2$. The endpoint must coincide with the
endpoint of $\mathbf u$.

\begin{center}
\begin{tikzpicture}[scale=0.95,>=stealth]
    \coordinate (O) at (0,0);
    \coordinate (B1) at (2.0,0.8);
    \coordinate (B2) at (0.8,1.8);
    \coordinate (TwoB1) at (4.0,1.6);
    \coordinate (U) at (4.8,3.4);

    % Oblique grid
    \foreach \i in {-1,0,1,2,3} {
        \draw[gray!45,dashed]
            ({-0.8+2.0*\i},{-1.8+0.8*\i}) --
            ({2.4+2.0*\i},{5.4+0.8*\i});
    }
    \foreach \j in {-1,0,1,2,3} {
        \draw[gray!45,dashed]
            ({-2.0+0.8*\j},{-0.8+1.8*\j}) --
            ({6.0+0.8*\j},{2.4+1.8*\j});
    }

    \draw[blue!70!black,line width=1.1pt,->] (O) -- (B1)
        node[midway,below] {$\mathbf b_1$};
    \draw[orange!85!black,line width=1.1pt,->] (O) -- (B2)
        node[midway,left] {$\mathbf b_2$};

    \draw[blue!70!black,line width=1.1pt,->] (O) -- (TwoB1)
        node[midway,below right] {$2\mathbf b_1$};
    \draw[orange!85!black,line width=1.1pt,->] (TwoB1) -- (U)
        node[midway,right] {$\mathbf b_2$};
    \draw[very thick,->] (O) -- (U)
        node[midway,above left] {$\mathbf u$};

    \draw[blue!70!black,dashed] (B2) -- (U);
    \draw[orange!85!black,dashed] (TwoB1) -- (U);

    \fill (O) circle (2pt) node[below left] {$O$};
    \fill (U) circle (2pt);
    \node[align=left,anchor=west] at (5.3,3.0)
        {$\mathbf u=2\mathbf b_1+\mathbf b_2$\\
         $[\mathbf u]_{\mathcal B}=(2,1)$};
\end{tikzpicture}
\end{center}

The coefficients are read from the oblique parallelogram. They are not lengths
of perpendicular shadows. They tell us how many signed copies of each basis
vector are needed.

\subsubsection*{The same vector on two different grids}

Let the ordinary Cartesian basis be
\[
\mathcal E=(\mathbf e_1,\mathbf e_2),
\qquad
\mathbf e_1=(1,0),
\qquad
\mathbf e_2=(0,1).
\]
Choose also the slanted basis
\[
\mathcal B=(\mathbf b_1,\mathbf b_2),
\qquad
\mathbf b_1=(2,1),
\qquad
\mathbf b_2=(1,2).
\]
Consider the geometric vector
\[
\mathbf u=(5,4)
\]
when it is described in the Cartesian basis.

On the rectangular grid, the endpoint is found by moving five units in the
$\mathbf e_1$ direction and four units in the $\mathbf e_2$ direction. Hence
\[
[\mathbf u]_{\mathcal E}=(5,4).
\]

On the slanted grid, we must find numbers $c_1,c_2$ such that
\[
\mathbf u=c_1\mathbf b_1+c_2\mathbf b_2.
\]
The two pictures below show the same arrow. Only the grid used to read its
coordinates has changed.

\begin{center}
\begin{tikzpicture}[scale=0.72,>=stealth]
    % Left: Cartesian grid
    \begin{scope}[xshift=0cm]
        \draw[gray!35,step=1] (-0.3,-0.3) grid (5.7,4.7);
        \draw[axis,->] (-0.3,0) -- (5.9,0) node[right] {$\mathbf e_1$};
        \draw[axis,->] (0,-0.3) -- (0,4.9) node[above] {$\mathbf e_2$};
        \coordinate (OL) at (0,0);
        \coordinate (UL) at (5,4);
        \draw[helper,dashed] (UL) -- (5,0);
        \draw[helper,dashed] (UL) -- (0,4);
        \draw[very thick,->] (OL) -- (UL)
            node[midway,above left] {$\mathbf u$};
        \node[below] at (5,0) {$5$};
        \node[left] at (0,4) {$4$};
        \node[align=center] at (2.7,-1.0)
            {rectangular grid\\$[\mathbf u]_{\mathcal E}=(5,4)$};
    \end{scope}

    % Right: oblique grid
    \begin{scope}[xshift=10.2cm]
        \coordinate (OR) at (0,0);
        \coordinate (B1) at (2,1);
        \coordinate (B2) at (1,2);
        \coordinate (TwoB1) at (4,2);
        \coordinate (UR) at (5,4);

        \foreach \i in {-1,0,1,2,3} {
            \draw[gray!45,dashed]
                ({-1+2*\i},{-2+\i}) -- ({3+2*\i},{6+\i});
        }
        \foreach \j in {-1,0,1,2,3} {
            \draw[gray!45,dashed]
                ({-2+\j},{-1+2*\j}) -- ({6+\j},{3+2*\j});
        }

        \draw[blue!70!black,line width=1pt,->] (OR) -- (B1)
            node[midway,below right] {$\mathbf b_1$};
        \draw[orange!85!black,line width=1pt,->] (OR) -- (B2)
            node[midway,left] {$\mathbf b_2$};
        \draw[blue!70!black,line width=1pt,->] (OR) -- (TwoB1)
            node[midway,below right] {$2\mathbf b_1$};
        \draw[orange!85!black,line width=1pt,->] (TwoB1) -- (UR)
            node[midway,right] {$\mathbf b_2$};
        \draw[very thick,->] (OR) -- (UR)
            node[midway,above left] {$\mathbf u$};

        % Parallel coordinate readings
        \draw[orange!85!black,dashed] (UR) -- (TwoB1);
        \draw[blue!70!black,dashed] (UR) -- (B2);

        \node[align=center] at (2.8,-1.0)
            {slanted grid\\$[\mathbf u]_{\mathcal B}=(2,1)$};
    \end{scope}
\end{tikzpicture}
\end{center}

On the rectangular grid we read the coordinates by projecting parallel to the
coordinate axes. On the slanted grid we do the analogous construction, but the
lines through the endpoint are drawn parallel to the other basis vector. This
produces an oblique parallelogram rather than a rectangle.

\subsubsection*{Finding the new coordinates by a system of equations}

The geometric picture tells us what must be true, but it does not yet calculate
$c_1$ and $c_2$. Write the basis vectors in ordinary Cartesian coordinates:
\[
\mathbf b_1=(a_1,a_2),
\qquad
\mathbf b_2=(b_1,b_2),
\qquad
\mathbf u=(x,y).
\]
The equation
\[
\mathbf u=c_1\mathbf b_1+c_2\mathbf b_2
\]
means equality of the two Cartesian components. Therefore
\[
\boxed{
\begin{aligned}
 a_1c_1+b_1c_2&=x,\\
 a_2c_1+b_2c_2&=y.
\end{aligned}}
\]
This ordinary system of two linear equations is the complete algebraic content
of changing from the Cartesian basis to the basis $\mathcal B$.

For our example,
\[
\mathbf b_1=(2,1),
\qquad
\mathbf b_2=(1,2),
\qquad
\mathbf u=(5,4),
\]
so
\[
\begin{aligned}
2c_1+c_2&=5,\\
c_1+2c_2&=4.
\end{aligned}
\]
From the first equation,
\[
c_2=5-2c_1.
\]
Substituting into the second gives
\[
c_1+2(5-2c_1)=4,
\]
so
\[
-3c_1=-6,
\qquad
c_1=2.
\]
Then
\[
c_2=5-2\cdot2=1.
\]
Thus
\[
\boxed{[\mathbf u]_{\mathcal B}=(2,1)}.
\]
The calculation agrees with the oblique parallelogram in the figure.

\subsubsection*{A general algebraic formula in the plane}

The same system can be solved symbolically. Starting from
\[
\begin{aligned}
 a_1c_1+b_1c_2&=x,\\
 a_2c_1+b_2c_2&=y,
\end{aligned}
\]
multiply the first equation by $b_2$ and the second by $b_1$:
\[
\begin{aligned}
 a_1b_2c_1+b_1b_2c_2&=xb_2,\\
 a_2b_1c_1+b_1b_2c_2&=yb_1.
\end{aligned}
\]
Subtracting eliminates $c_2$:
\[
(a_1b_2-a_2b_1)c_1=xb_2-yb_1.
\]
Similarly, eliminating $c_1$ gives
\[
(a_1b_2-a_2b_1)c_2=a_1y-a_2x.
\]
Hence
\[
\boxed{
 c_1=\frac{xb_2-yb_1}{a_1b_2-a_2b_1},
 \qquad
 c_2=\frac{a_1y-a_2x}{a_1b_2-a_2b_1}
}.
\]

The denominator
\[
a_1b_2-a_2b_1
\]
is nonzero precisely because the basis vectors are not parallel. If it were
zero, the two chosen directions would collapse onto one line and the
coordinates would not be uniquely determined.

\subsubsection*{Changing directly from one nonstandard basis to another}

Suppose
\[
\mathcal B=(\mathbf b_1,\mathbf b_2)
\qquad\text{and}\qquad
\mathcal C=(\mathbf c_1,\mathbf c_2)
\]
are two bases. Assume that
\[
[\mathbf u]_{\mathcal B}=(p,q).
\]
Then first reconstruct the geometric vector:
\[
\mathbf u=p\mathbf b_1+q\mathbf b_2.
\]
If
\[
\mathbf b_1=(a_1,a_2),
\qquad
\mathbf b_2=(b_1,b_2),
\]
then
\[
\mathbf u=(pa_1+qb_1,\,pa_2+qb_2).
\]
Now seek numbers $r,s$ such that
\[
\mathbf u=r\mathbf c_1+s\mathbf c_2.
\]
Writing
\[
\mathbf c_1=(d_1,d_2),
\qquad
\mathbf c_2=(e_1,e_2),
\]
we solve
\[
\begin{aligned}
 d_1r+e_1s&=pa_1+qb_1,\\
 d_2r+e_2s&=pa_2+qb_2.
\end{aligned}
\]
The solution $(r,s)$ is exactly $[\mathbf u]_{\mathcal C}$.

Thus a change from one basis to another is not a mysterious new operation. It
is two familiar steps:
\[
\boxed{
\text{old coordinates}
\longrightarrow
\text{the geometric vector}
\longrightarrow
\text{new coordinates}.}
\]

\subsubsection*{Why perpendicular projection is special}

For the Cartesian basis, the coordinate directions are perpendicular unit
vectors. Therefore the horizontal and vertical coordinates can be read from
ordinary perpendicular projections.

For a slanted basis, this is no longer true. Projecting $\mathbf u$
perpendicularly onto $\mathbf b_1$ and $\mathbf b_2$ produces two shadows, but
those shadows generally do not add back to $\mathbf u$. The correct geometric
construction uses lines parallel to the other basis vector, exactly as in the
oblique grid above.

\begin{remarkbox}
Later, after matrices have been introduced, the same systems of equations will
be written much more compactly. Matrix notation will not create a new method;
it will only abbreviate the elimination calculations performed here.
\end{remarkbox}

\begin{summarybox}
A basis is a pair of nonparallel coordinate directions. The same geometric
vector may have different coordinate pairs in different bases. To find the
coordinates in a slanted basis, write
\[
\mathbf u=c_1\mathbf b_1+c_2\mathbf b_2
\]
and solve the resulting two scalar equations. Rectangular coordinates use
perpendicular projections; slanted coordinates use an oblique parallelogram
whose sides are parallel to the basis vectors.
\end{summarybox}
